% ============================== Part A ==============================
\problem{A1}{Two measurements of one number}

\begin{given}
A measurement $d$ of an unknown $x$ with Gaussian noise of standard deviation $\sigma$ has likelihood
\[
L(x) \propto \exp\!\left[-\frac{(x-d)^2}{2\sigma^2}\right].
\]
Independent measurements multiply: the joint likelihood of $d_1$ and $d_2$ is $L_1(x)\,L_2(x)$. Maximising it is the
same as minimising $\chi^2 = \sum_i (d_i - x)^2/\sigma_i^2$, the weighted least-squares problem.
\source{\cHBLwls.}
\end{given}

\need{nothing; this is the start.}

\begin{enumerate}
  \item Show that $L_1(x)L_2(x)$ is again a Gaussian in $x$. Find its mean $\hat x$ and variance $\sigma_{\hat x}^2$.
  \item Write $\hat x$ as a weighted average of $d_1$ and $d_2$. What are the weights?
  \item A KL galaxy measures one shear component as $d_1 = 0.03 \pm 0.05$; a second galaxy gives $d_2 = 0.01 \pm 0.02$.
    Compute $\hat x$ and $\sigma_{\hat x}$. Which galaxy dominates, and why?
\end{enumerate}

\solution{
(a) Adding the exponents,
$-\tfrac12\big[(x-d_1)^2/\sigma_1^2 + (x-d_2)^2/\sigma_2^2\big] = -\tfrac12 (w_1+w_2)\,x^2 + (w_1d_1 + w_2d_2)\,x + \text{const}$,
with $w_i = 1/\sigma_i^2$. Completing the square gives a Gaussian with
\[
\hat x = \frac{w_1 d_1 + w_2 d_2}{w_1 + w_2},\qquad \sigma_{\hat x}^2 = \frac{1}{w_1 + w_2}.
\]
(b) The weights are the precisions $w_i = 1/\sigma_i^2$, normalised to sum to 1. Precisions add.
(c) $w_1 = 400$, $w_2 = 2500$: $\hat x = (12 + 25)/2900 = 0.0128$ and $\sigma_{\hat x} = 2900^{-1/2} = 0.0186$. The second
galaxy dominates: its weight is $(0.05/0.02)^2 = 6.25$ times larger. The combined error is smaller than either.
}

% --------------------------------------------------------------------
\problem{A2}{Many measurements}

\begin{given}
The result of A1 generalises to $N$ independent Gaussian measurements $d_i \pm \sigma_i$ of one $x$: the product of
the likelihoods is Gaussian, with precision $\sum_i w_i$ and mean $\sum_i w_i d_i / \sum_i w_i$, where
$w_i = 1/\sigma_i^2$.
\source{\cHBLwls.}
\end{given}

\need{A1 (precisions add; the mean is precision-weighted).}

\begin{enumerate}
  \item Show that for equal errors $\sigma_i = \sigma$ the combined error is $\sigma/\sqrt N$.
  \item A KL galaxy measures each shear component with $\sigma_g \approx 0.05$. How many such galaxies are needed for
    an error of $0.01$? What is the error for a cell of 5 galaxies?
  \item Our cluster cells hold about 5 galaxies. The subhalo changes the shear in most cells by $\lesssim 0.01$. Can a
    single cell detect the subhalo?
\end{enumerate}

\solution{
(a) $\sum_i w_i = N/\sigma^2$, so $\sigma_{\hat x} = \sigma/\sqrt N$.
(b) $0.05/\sqrt N = 0.01$ gives $N = 25$. With 5 galaxies, $0.05/\sqrt5 = 0.022$. This is the per-cell error
seen in the cluster shear maps.
(c) No: a $\lesssim 0.01$ signal against a $0.022$ error is below $0.5\sigma$ per cell. Detecting the subhalo needs
the information of many cells (or many galaxies) combined. That is what the $\kappa$ map does (Parts E--G).
}

% --------------------------------------------------------------------
\problem{A3}{Two-component measurements with covariance matrices}

\begin{given}
A measurement $\vd_i$ of a 2-vector $\vx$ (for us, $\vx = (g_1, g_2)$) with Gaussian noise of covariance $C_i$ has
likelihood $L_i(\vx) \propto \exp\!\big[-\tfrac12 (\vx-\vd_i)^\T F_i\,(\vx - \vd_i)\big]$, with precision matrix
$F_i = C_i^{-1}$. The general weighted least-squares solution, for data $Y$ with covariance $C$ and a linear model
$Y = AX$, is $\hat X = [A^\T C^{-1} A]^{-1} A^\T C^{-1} Y$, with covariance $[A^\T C^{-1} A]^{-1}$.
\source{\cHBLwls.}
\end{given}

\need{A1, A2.}

\begin{enumerate}
  \item Show that $\prod_i L_i(\vx)$ is Gaussian with precision $\Ftot = \sum_i F_i$ and mean
    $\hat\vx = \Ftot^{-1}\sum_i F_i \vd_i$. (Either complete the square, or use the weighted least-squares solution
    with $A$ a stack of $2\times2$ identity matrices.)
  \item Two galaxies: $\vd_1 = (0.02, 0)$ with $F_1 = \mathrm{diag}(2500, 100)$; $\vd_2 = (0, 0.02)$ with
    $F_2 = \mathrm{diag}(100, 2500)$. Compute $\hat\vx$. Compare with a ``scalar'' average that weights each galaxy by
    the trace of its $F_i$.
  \item Now $\vd_1 = (0.02, 0)$ with $F_1 = \begin{pmatrix}1300 & 1200\\ 1200 & 1300\end{pmatrix}$ and
    $\vd_2 = (0,0)$ with $F_2 = 100\,I$. Compute $\hat\vx$. Neither measurement has a non-zero second component; why
    is $\hat x_2 \neq 0$?
  \item KL kinematics measure one shear direction much better than the other. What does this problem say about
    combining KL galaxies?
\end{enumerate}

\solution{
(a) $\sum_i (\vx - \vd_i)^\T F_i (\vx - \vd_i) = \vx^\T \Ftot \vx - 2\vx^\T \sum_i F_i\vd_i + \text{const}
= (\vx - \hat\vx)^\T \Ftot (\vx - \hat\vx) + \text{const}$, with $\hat\vx = \Ftot^{-1}\sum_i F_i\vd_i$.
(b) $\Ftot = \mathrm{diag}(2600, 2600)$ and $\sum F_i\vd_i = (50, 50)$, so $\hat\vx = (0.0192, 0.0192)$. Each galaxy
dominates the component it measures well. The trace weights are equal (2600 each), giving $(0.01, 0.01)$: half the
signal is lost, because the scalar average throws away which direction each galaxy measures.
(c) $\Ftot = \begin{pmatrix}1400 & 1200\\1200 & 1400\end{pmatrix}$ and $\sum F_i\vd_i = (26, 24)$. Then
$\Ftot^{-1} = \frac{1}{520000}\begin{pmatrix}1400 & -1200\\ -1200 & 1400\end{pmatrix}$, so
$\hat\vx = (0.0146, 0.0046)$. $F_1$ has off-diagonal terms: galaxy 1 measures the combination $x_1 + x_2$ well
(eigenvalue 2500 along $(1,1)/\sqrt2$) and $x_1 - x_2$ poorly (eigenvalue 100). The components mix.
(d) Matrix weights are required; scalar weights mislabel the result. The same matrix weighting reappears in the
reference shear (B4) and in BA14 (G1).
}

% --------------------------------------------------------------------
\problem{A4}{Posterior, prior and likelihood}

\begin{given}
Bayes' rule: $p(x \mid d) = L(x)\,\pi(x) / Z$, with $L(x) = p(d\mid x)$ the likelihood, $\pi(x)$ the prior and $Z$
a normalising constant. A neural posterior estimator (NPE) trained on simulations whose parameters were drawn from
$\pitrain$ learns $q(x\mid d) \approx p(x\mid d)$ under that prior. If posterior samples were drawn under an
\emph{interim} prior, dividing by that prior (reweighting the samples) recovers the likelihood.
\source{\cHoggBayes; \cInterim.}
\end{given}

\need{nothing beyond the definitions.}

\begin{enumerate}
  \item Show that $q(x\mid d)/\pitrain(x) \propto L(x)$ wherever $\pitrain(x) > 0$.
  \item Our training prior is uniform on the box $|g_k| < 0.1$. What does $q/\pitrain$ look like inside the box, and
    what do we know about $L$ outside it?
  \item Why must we divide by the prior the network was \emph{trained} with, not by some other prior we prefer?
\end{enumerate}

\solution{
(a) $q = L\,\pitrain/Z$, so $q/\pitrain = L/Z \propto L$ (constant in $x$).
(b) Inside the box $\pitrain$ is constant, so $q/\pitrain \propto q$: the likelihood has the posterior's shape there.
Outside the box $q = 0$, so the likelihood is unknown there: the NPE says nothing about shears it never saw.
(c) The network has absorbed $\pitrain$ into $q$. Dividing by a different $\pi'$ leaves a factor $\pitrain/\pi'$
that is not part of the likelihood, which would bias every later product of likelihoods.
}

% --------------------------------------------------------------------
\problem{A5}{\py\ Why one galaxy's posterior mean is shrunk}

\begin{given}
Take one shear component. The likelihood is Gaussian, $L(x) \propto \exp[-(x-d)^2/2\sigma^2]$, with $\sigma = 0.07$
(the untruncated width of a typical KL likelihood), and the prior is uniform on $|x| < 0.1$. The posterior is then a
Gaussian truncated to $[-0.1, 0.1]$. Its mean is the ``posterior-mean shear estimate'' of one galaxy. The data scatter
as $d \sim \mathcal N(x_{\rm true}, \sigma^2)$. \emph{Derived here}.
\end{given}

\need{A4 (posterior = likelihood $\times$ prior).}

\begin{enumerate}
  \item Let $\alpha = (-0.1 - d)/\sigma$, $\beta = (0.1 - d)/\sigma$, and let $\phi$ and $\Phi$ be the standard normal
    density and CDF. Show that the posterior mean is
    $d + \sigma\,[\phi(\alpha) - \phi(\beta)]/[\Phi(\beta) - \Phi(\alpha)]$. (Integrate $x\,e^{-(x-d)^2/2\sigma^2}$ by
    substitution.)
  \item Explain without calculation why its average over $d$ is smaller in magnitude than $x_{\rm true}$.
  \item \py\ Compute $\E_d[\text{posterior mean}]$ for $x_{\rm true} \in [-0.05, 0.05]$, fit a line
    $(1+m)x_{\rm true} + c$, and report $m$.
  \item \py\ Now average $N$ independent galaxies at the same $x_{\rm true}$ by multiplying their likelihoods, which is
    a Gaussian of width $\sigma/\sqrt N$ by A2, truncated to the box. Repeat (c) for $N = 1, 2, 4, 8, 16$. How large
    must $N$ be for $|m| < 0.01$?
\end{enumerate}

\solution{
(a) With $u = (x-d)/\sigma$, $\int_{-0.1}^{0.1} x\,\phi(u)\,dx \propto d\,[\Phi(\beta)-\Phi(\alpha)] +
\sigma\int_\alpha^\beta u\,\phi(u)\,du$, and $\int_\alpha^\beta u\phi\,du = \phi(\alpha) - \phi(\beta)$. Dividing by the
normalisation $\Phi(\beta) - \Phi(\alpha)$ gives the result.
(b) When $d$ falls near or beyond an edge, the box cuts off the far tail of the likelihood, so the mean is dragged
inward. For $x_{\rm true} > 0$ the upper edge is reached more often than the lower one, so on average the estimate
is pulled toward zero.
(c) $m \approx -0.57$ for $\sigma = 0.07$. A single posterior mean is strongly shrunk, close to the $m \approx -0.6$
of the raw NPE posterior means (paper Sec.~6.1).
(d) $m \approx -0.57, -0.36, -0.17, -0.052, -0.0085$ for $N = 1, 2, 4, 8, 16$ (and $-0.0005$ for $N = 32$). The product
narrows as $\sigma/\sqrt N$ and stops touching the box, so the shrinkage disappears. With every galaxy at
$\sigma = 0.07$, $|m| < 0.01$ needs $N \approx 16$. Real KL widths vary from galaxy to galaxy, and the narrow ones
dominate a product, which is why the paper's stacks reach percent-level $m$ at smaller $N$. Multiplying likelihoods,
rather than averaging posterior means, is what removes the bias.
}
